\cotefile{161-2}
\scanpage{20}
% archive sheet number: 19

% L1–3
$1^\circ)$ \deleted{\unsure{Théorie vi}} Structure vide sur objets de base $(X_i)_{i\in I}$. \unsure{Représent.ble|Représentables} par un topos $R_0=$ ? Pb 1 ; on a $R_0=\mathrm{Ens.}$ \deleted{\unsure{Ens}} $I$-multi\inserted{\unsure{co|$\omega$}}simpliciaux \inserted{\unsure{opposé-dual}} $={}\in\widehat{\mathbf J}_0$, $\mathbf J_0=$ catégorie \deleted{des} \inserted{familles d'}ensembles finis $(E_i)_{i\in I}$ \unsure{presque tous vide}.
% L2–4, left (slanted)
% slanted
\margin{$J=$ catégorie stricte \unsure{t.q.} mais \unsure{seulemt équivalence} pour $I$}
% L3–4, interline
\[
 I\to\mathbf J_0\to R_0\mathrel{:=}\widehat I\longrightarrow R
\]
$\mathbf J$ est \uline{générateur dans} $R$. \quad $I-\underset{\smile}{\overset{R_0}{J_0}}$

% L4–5
$2^\circ)$ On regarde des flèches \deleted{\unsure{entre des}} entre objets de $R_0$ provenant de $\mathbf J_0$ (celles de $R$), d'où \deleted{$R_{\unsure{1|2|0}}$} \inserted{$J_1$ et $R_1=\rho_\Delta(J_1)$} \deleted{$I\unsure{\to}J_0\unsure{\to}R_0\to R_1\to R$}

% L6–7
$3^\circ)$ on met des axiomes d'égalité sur \deleted{des} \inserted{les} composés (en vertu de $R$) d'où
\[
 \mathbf J\to R_0\to R_1-R_2\to R
\]

% L8
\[
 T_{R_1}(E)=\underline{\mathrm{Hom}}_{\unsure{B|\beta|\Delta}}(R_1,E)\simeq
 \underline{\mathrm{Hom}}(\mathbf J_1,R)\hookleftarrow
 \underline{\mathrm{Hom}}_{\Delta}(R,E)=T_R(E)
\]

% L9–10
où $\mathbf J_1$ est la sous-catégorie pleine de $R$ définie par $b:\mathbf I\to R$, dont $\mathbf R_{2}$ est la \unsure{Ds-enveloppe|$\Delta$-enveloppe} de la petite catégorie $\mathbf J_1$ (cf Pb.1)

% L11–12
\deleted{\unsure{$4^\circ$}} Maintenant \deleted{\unsure{la transformation}} \inserted{le foncteur} $T_R\longrightarrow T_{R_2}$ est pl. fidèle, on sait donc par le N° précédent qu'on a fini par

% L13–20
$4^\circ)$ On met d'autres axiomes \deleted{\illegible{1}} sur la structure. [on peut expliciter lesquels : a) foncteurs \uline{continus} $I_1\unsure{-|\to}E$, i.e. \deleted{satisf} qui transforment certains \deleted{systèmes} cônes inductifs de $I_1$ \inserted{(ceux existant dans $R$)} en cônes exacts \inserted{de $E$} ; \deleted{b) foncteurs} ce sont les foncteurs qui se prolongent en des foncteurs \uline{commutant aux} \uline{lims} $R\to E$ ; plus la condition de transformer certains \deleted{diagrammes de} cônes projectifs finis \uline{de $E$} (ou si on veut, de $R_2'$, le topos cl. \unsure{défini|déf.} par ce système d'axiomes) en cônes exacts ..... . On trouve $I-R_0-R_1-R_2-R_3\simeq R$

% L11–22, left (sideways)
% margin note written sideways
\margin{\diagram{\begin{tikzcd}[ampersand replacement=\&,column sep=small,row sep=small]
 \rho(I) \arrow[d,equal] \& \rho(J_1\unsure{;}J_0) \arrow[d,equal] \& \rho(J_2)\arrow[d,equal] \& {}\\
 R_0\arrow[r] \& R_1\arrow[r,dash] \& R_2\arrow[r,dash]\arrow[d] \& R_3=R\arrow[dl]\\
 J_0\arrow[r]\arrow[u,leftrightarrow] \& J_1\arrow[r]\arrow[u]\arrow[ur,phantom,"\unsure{J_1'}" description]\& J_2 \& {}\\
 I\arrow[u,dash] \& {} \& {} \& {}
 \end{tikzcd}}\deleted{\illegible{1}} $\Delta$-enveloppe, \unsure{pour toutes|pour toutes les} sommes \unsure{provenant de|provenant des} $J_0)=\rho(J'_1)$ $J'_1$ enveloppe \deleted{\unsure{engend|engendré}} par $J_0$ et les flèches données.}

% L21–27
\uline{NB} Comme on a rajouté dans $3^\circ$ et $4^\circ$ des axiomes, on peut les condenser en une seul pas, grâce au fait que $R_1-R_2$ est surjectif sur objets et flèches (cf. Pb 2). Donc on peut \inserted{tout} faire par les trois pas minimum : objets, se \unsure{prendre} flèches \inserted{entre produits finis}, axiomes sur les flèches, exprimables en termes \deleted{de lims} d'inversion de flèches entre \deleted{\unsure{entre}} \uuline{\unsure{paires|termes}} \unsure{formés|formées} entre ces produits \deleted{lims \illegible{3}} \inserted{pour des syst. inductifs de} ou morphismes entre tels formés grâce aux flèches introduites dans les données.

% L25–27, lower left
\margin{\uline{NB} $\widehat{J_1}\to R$ est un \uline{foncteur} de \uline{localisation}}
\scanpage{21}

% L1
\deleted{\uline{\unsure{Nécessité}} de \unsure{faire}}

% L2–9
\uline{Re NB} Si on veut seulement montrer que \deleted{les} topos $R$ est ``algébrique'' c'est encore plus court : On choisit $C$ petite avec \uline{lims} finies, telle que $R\simeq\widetilde C$, \deleted{et} on \deleted{écrit que} part de la structure algébrique $T_C^{\Delta{}_0}$ -- sa restriction \unsure{,} $\Delta$ est donc algébrique, elle est représentée par $\widehat C$, cf \deleted{\unsure{P}} n° 1 -- puis on rajoute des axiomes exprimant la \uline{contin}uité d'un foncteur $C-E$ (commutant déjà aux \uline{lims} finies) par le fait de transformer familles couvrantes en \unsure{famille|familles} couvrantes ...

% L10
\deleted{Je \unsure{note à présent en résumé} que l'algébricité}

% L11
\deleted{\uline{Nécessité}} Il faut suivre pas à pas la construction transfinie

% L12–14
a) Structure vide sur ens de base \inserted{$I$}. Représentée par $\widehat{\widetilde I}$ (cf Pb 1), \inserted{où $\widetilde I=\rho_{\Delta_0}(I)$.} \deleted{$I$ engendre $\Delta$-engendre} \uline{NB} $I$ $\Delta$-engendre $\widehat{\widetilde I}$ en deux pas : rajouter des produits, puis des \uline{lim} quelc.

% L15–17
b) \deleted{\unsure{Ra}} Mettre axiomes sur théorie algébrique $T_R$, $R$ un topos, cela relève du Pb 2 traité. On trouve $R'$ et $R-R'$ \unsure{ess.|en} surjectif (et aussi surjectif sur les Hom).

% L18–22
c) Rajouter flèches (petit ensemble \inserted{$H$}) : \deleted{\unsure{Re.|$R$}} théorie algébrique $T_R$, $R$ \deleted{déjà} topos. \deleted{Soit $C$ l'ens des} Soit \deleted{$C\subset R$ la petite sous-catég} $\mathbf I$ \inserted{$\subset\operatorname{Ob}R$} l'ens des sources et buts de ces flèches, $C$ la catégorie libre engendrée par \deleted{les ob} $I$ comme objets et les flèches $H$, $\rho_\Delta(C)$ la catégorie enveloppante, alors la théorie $T'$ est le produit 2-fibré
% L23
\[
 T'\simeq T_R\mathop{\times}^{2}_{\tau_{\mathbf I}}\tau_C
 \qquad \deleted{\unsure{\tau_I=\underline{\mathrm{Hom}}_{\mathrm{Cat}}}}
\]
% L24–26
où $\tau_C(E)\simeq\underline{\mathrm{Hom}}_{\mathrm{Cat}}(C,E)$, et $\tau_{\mathbf J}(E)$ se définit de $\overline{\unsure{\mathrm m}}$, $J$ considéré comme catégorie \unsure{discrèt.|discrète}. Or (pb 1) $\tau_C\simeq T_{R_1}$, $\tau_I\simeq T_S$ $R_1$ et $S$ deux topos, donc il faut prouver l'existence
% archive sheet number: 20
\scanpage{22}

% L1–4
du produit 2-fibré dans (Top) $[\,\deleted{\unsure{de|des}}\,R\mathop{\times}^{2}_{S}R_1\,]$, ou la 2-somme amalgamée dans $\Delta$ \quad $R\amalg_S R_1$. Cette construction est $\pm$ connue, redonnons en une autre, dans le cas général $R\mathop{\times}^{2}_{T}S$, avec les foncteurs images inverses
% L4, right
\diagram{\begin{tikzcd}[ampersand replacement=\&,column sep=small,row sep=tiny]
 {} \& R\\ T\arrow[ur,"f^*"]\arrow[dr,"g^*"'] \& {}\\ {} \& S
\end{tikzcd}.}
% L4–7
On choisit des \deleted{familles} \inserted{petites sous-catégories} génératrices $T_0,R_0,S_0$ telles que $f^*(T_0)\subset R_0$, $g^{\unsure{*|'}}(T_0)\subset S_0$. On prend $U_0$ la catégorie-somme amalgamée de $R$ et $S$, de sorte que pour tte \inserted{catégorie} $E$ \deleted{\illegible{1}} \ldots .
% L8
\[
 \underline{\mathrm{Hom}}(U_0,E)\xrightarrow{\sim}
 \underline{\mathrm{Hom}}_{\deleted{\unsure{\Delta}}}(R_0,E)
 \mathop{\times}^{2}_{\underline{\mathrm{Hom}}(T_0,E)}
 \underline{\mathrm{Hom}}(S_0,E)
\]

% L9–14
Il faut, si $E\in\operatorname{Ob}\Delta$, exprimer que \unsure{l'un|un} objet du \deleted{second \unsure{membre|mb.}} premier mb, $U_0\xrightarrow{h}E$, est dans \deleted{$\underline{\mathrm{Hom}}_\Delta(R,E)$} la sous-catégorie pleine $\underline{\mathrm{Hom}}_\Delta(R,E)\mathop{\times}^{2}_{\underline{\mathrm{Hom}}_\Delta(T,E)}\underline{\mathrm{Hom}}_\Delta(\unsure{T|S},E)$ du deuxième mb. Si on a pris $R_0,S_0$ stables par \uline{lim} finies, cela s'exprime aisément par la \deleted{commutation de $h$ aux} condition que $h$ transforme certains cônes projectifs finis de $U_0$ (\inserted{\unsure{ceux|aux}} provenant des cônes exacts \deleted{de} $R_0$ et de $S_0$) en cônes exacts,
% L15–19
et \deleted{que les \unsure{im|uns}} qu'elle transforme certaines familles (provenant de familles couvrantes de $R_0$ et $S_0$) en familles couvrantes. Donc on réalise la théorie envisagée $\tau=T_R\mathop{\times}^{(2)}_{T_T}T_S$ comme sous-théorie pleine (définie par des axiomes d'\deleted{\unsure{égalité}}isomorphie, finitaires) de $T_{\rho_\Delta(U_0)}$, ce qui donne la conclusion.

% L19–26
Mais sachant désormais que \deleted{\unsure{Dans|dans}} la théorie $\tau$ est \unsure{représent.ble|représentable} par un topos, \deleted{et} revenant à la description générale comme sous-théorie d'un $T_{\rho_\Delta(U_0)}$ (sans \inserted{plus de} condition que $R_0$, $S_0$ stables par $\underleftarrow{\lim}$ existantes, on voit que c'est donné par un \uline{sous}-\uline{topos} de $\rho_\Delta(U_0)$. Ainsi on trouve, si $T_0\to S_0$ est surjectif sur les objets (comme dans le cas qui nous intéressait) que $R_0-U_0$ est surjectif sur les objets, donc $R\to U=R'$ \deleted{est géné} a une image qui est génératrice.
% archive sheet number: 21
\scanpage{23}

% L1–6
d) \uline{Passage à la limite} Cela correspond encore : des 2-Limites \inserted{inductives} \unsure{dans|\deleted{dans}} $\Delta$, il faut montrer que les 2-limites transfinies existent et sont des topos. La démonstration donnée dans c, marche en fait pour tt système inductif \deleted{de topos} dans $\Delta$, formé de topos [du système projectif dans (Top)] : $(R_\alpha)$ : On se débrouille pour choisir des
% L7–11
\deleted{systèmes de} \inserted{petites cat.} généra\deleted{trices} \inserted{$C_\alpha\subset R_\alpha$} \unsure{qui|\deleted{qui}} sont stables par les morphismes de transition \deleted{\unsure{ce qui est f}} -- p.ex. à l'aide des \unsure{fibrations|filtrations} \unsure{coreliantes|croissantes|cardinales}, mais c'est un marteau pour une mouche -- et on prend \deleted{la catégorie-\unsure{limite in|2-} de} $C=\underset{\alpha\,(\mathrm{Cat})}{\underline{\mathrm{Lim}}}\, C_\alpha$, on \deleted{prend} trouve un sous-topos de $\widehat C$.

% L12–17
En l'occurrence, \deleted{on voit que si on} \inserted{dans le cas filtrant} on voit que si \deleted{un} $R$ est le topos limite, \deleted{et si} les images dans $R$ des $R_\alpha$ forment un syst. de générateurs ; \deleted{\unsure{hyp si}} \inserted{\unsure{donc|alors|d'où}} on a une \deleted{syst. de générateurs} catégorie $\mathbf J$ s'envoyant dans le syst. inductif des $R_\alpha$, $J-R_\alpha$ engendre $R_\alpha$ pour tt $\alpha$, alors $J-R$ engendre $R$.

% L18, horizontal rule
\par\noindent\rule{\linewidth}{.3pt}\par

% L18–24, lower diagram
\diagram{\begin{tikzcd}[ampersand replacement=\&,column sep=small,row sep=small]
 \rho(I)=\widehat J_0\arrow[d,equal] \& \rho'(J_1)\arrow[d,equal]\& \rho^{\prime}(J_2)\arrow[d,equal]\& {}\\
 R_0\arrow[r] \& R_1\arrow[r,"\mathrm{loc}"]\& R_2\arrow[r,"\mathrm{loc}"]\& R_3=R\\
 J_0\arrow[u]\arrow[r,"\substack{\text{bij sur}\\\text{objets}}"]\&
 J_1\arrow[u]\arrow[r,"\mathrm{loc.}","{(\text{bij sur obj})}"']\&
 J_2\arrow[u]\arrow[ur,leftrightarrow]\& {}\\
 I\arrow[u]\& {} \& {} \& {}
\end{tikzcd}}

% L18, annotation above diagram
$\rho'$ la $\Delta_\pi$-\unsure{b|\Delta}-enveloppe, où $\Delta_\pi$ cat. avec produits finis

% L21–25, beneath J1 and J2
$J_1$ $=$ analyseur engendré par $J_0$ (analyseur \unsure{vierge|\deleted{libre}}) et les flèches données -- analyseur libre
\qquad $J_2$ $=$ sous-analyseur plein de $R$ engendré par $I\to R$, ens. sous-jacent $J_0$

% L19–25, lower right
\margin{$R$ apparait comme une catégorie de foncteurs de \deleted{\unsure{préfaisc}} $\widehat J_1$, où $J_1$ est l'analyseur libre sur $I$ engendré par les flèches induites (sur $J_0$) par $R$, $\widehat J_1$}

% L26–27, foot
\diagram{\begin{tikzcd}[ampersand replacement=\&,row sep=small,column sep=small]
 \widehat J_1\arrow[r,dash,"\mathrm{loc?}"]\arrow[d,"\mathrm{loc}"']\&\widehat J_2\arrow[d,"\mathrm{loc}"]\arrow[r,"\deleted{\unsure{\mathrm{loc}}}"]\&R\\
 R_1\arrow[r,dash,"\mathrm{loc}"]\&R_2\arrow[r,"\mathrm{loc}"]\&R_3
\end{tikzcd}}
% archive sheet number: 22
\scanpage{24}

% L1
\uline{Structure représentée par} $B_G=\widehat C\qquad C=\deleted{\unsure{(G,}}(\mathord{.},G)$

% L2
Relativisation dans un topos.

% L3
\uline{Topos produit}.

% L4–8
Remarques sur les $T(E)\simeq\underline{\mathrm{Hom}}\,\mathrm{top}(E,R)^o\simeq\underline{\mathrm{Hom}}_\Delta(R,E)$ : ces catégories \uline{sont stables par \uline{lim} filtrantes} (ça \unsure{t'élimine|élimine|il vient} : pas mal de conditions sur des structures : elles ne correspondent pas à des théories algébriques sur \deleted{quel que ce soit} \inserted{quelque} $\Delta'\supset\Delta$.

% L8, right
\uline{NB} En général aucun autre type de lim.

% L9–11
\uline{Question} Si $\forall E$, $T(E)$ admet des \deleted{produits fibrés quelconques} \inserted{$\underleftarrow{\lim}$ finies}, \inserted{(et les $T_R(E)\longrightarrow T_R(E')$ y commutent)} aura-t-on $T=$ restriction à $\Delta$ d'une théorie algébrique sur $\Delta_0$, i.e. $R\simeq\widehat C$, avec $C$ \inserted{(petite)} $\in\operatorname{Ob}\Delta_0$ i.e. $C$ stable par $\underleftarrow{\lim}$ finies ?

% L12–14
Si $T$ est de ce type, on sait que $T(E)$ est \unsure{$\simeq$} stable par $\underrightarrow{\lim}$ \deleted{\unsure{fini|finies}} quelconques et accessible, \deleted{(\uline{NB} Mor$(\unsure{C})$} et (donc) stable par $\underleftarrow{\lim}$ quelc.

% L15
$T_R(\mathrm{Ens})\simeq\underline{\mathrm{Pt}}(R)^o\simeq\underline{\mathrm{Fib}}(R)$. Peut être vide et $R\not\simeq\operatorname{Top}(\varnothing)$

% L16–19
\uline{Sous topos du topos libre \inserted{$R$} engendré par un objet} $X$ (espèces de structures vides sur un objet de base) : \deleted{\unsure{comme|comparer}} $R=\widehat C$, $C=$ catégorie opposée des ens finis, ou $R=$ ens. semi-simpliciaux. \inserted{\uline{Pts} : les Ens} Sous-Topos \inserted{ouverts} $=$ cribles de $\mathbf C$, on trouve les sous-topos suivants (ils sont caractérisés par leurs points) :

% L19–24, table
\diagram{$\begin{array}{r|lll}
 & \unsure{\mathbb R|\mathbf R|R}, & \varnothing, & \mathbb R^*=\unsure{\operatorname{Im}|I\mathrm m}(X\to e)\\
 \substack{\text{pts}\\\text{du p.f.}\\\text{(crible}}\ \}\to &
 \mathbb N, & \varnothing, & \mathbb N^*\\
 &\mathbf S, & \varnothing, & (\mathrm{Ens}_{f}\ \text{non vides})^*\\
 \text{Sous-théorie} & \unsure{T(E)=\mathrm{Ens}|T(E)=E|\deleted{T(E)=E}}, & \varnothing, & T(E)\subset E\\
 &&&\substack{\text{ss catégorie pleine}\\\text{des objets couvrant }E}
 \end{array}$}

% L21, left
\margin{$S=(\mathrm{Ens}_{f})^o$}

% L24–25, arrow below empty column
\diagram{\begin{tikzcd}[ampersand replacement=\&,row sep=tiny]
 {}\arrow[d] \\ \substack{T(E)=\varnothing\ \text{sauf}\\\text{si } E\simeq\mathrm{Top}(\varnothing)}
\end{tikzcd}}

% L26–28
\uline{Sous-topos fermés complémentaires}
\diagram{$\begin{array}{r|lll}
 &\varnothing,&\mathbb R,&\unsure{\mathbf R_0|\Phi_0}=\unsure{\complement|f|\mathcal C}U\\
 \text{p.p.f.}\Rightarrow &\varnothing,&\mathbb N,&\{\Delta_0\}\\
 &&&T(E)\subset E\\
 &&&\substack{\text{catég. ponctuelle}\\\text{réduite à }\deleted{\unsure{\text{l'obj}}}\text{l'objet vide}}
 \end{array}$}
% archive sheet number: 23
\scanpage{25}

% L1–4
Il y a enfin un sous topos $V$ \deleted{$\unsure{\Delta_1}$} qui n'est ni ouvert, ni fermé, \deleted{mais loc.} dont les pts sont $\Delta_0,\Delta_1$, et qui \deleted{\unsure{est}} est donné par
\[
 \begin{aligned}
 T_V(E)&=\text{cat. des sous objets de }\deleted{\unsure{X}}e_E\\
 &\text{(i.e. des objets }X\in\operatorname{Ob}E\text{ tels que }x\to x_e\text{ mono)}
 \end{aligned}
\]
% L5–7
et le sous-topos $\unsure{R^*|U^*|\mathrm{Ens}^{+}}\cap V$ \deleted{$\simeq\unsure{\varphi}$} topos ponctuel, point $\Delta_1$, tel que
\[
 \underset{\substack{\text{objets}\\\text{finaux}}}{T_{U\cap V}(E)}\subset E\qquad
 \text{i.e. }x\xrightarrow{\sim}e_E
\]

% L8
En résumé, on a $6$ sous-topos, avec diagramme d'implications
% L9–14, left diagram
\diagram{\begin{tikzcd}[ampersand replacement=\&,row sep=small,column sep=small]
 {} \& R\arrow[dl,dash]\arrow[dr,dash]\&{} \&{}\\
 R^*\arrow[dr,dash]\&{}\& V=\overline{(\varphi_0,\varphi_1)}\arrow[dl,dash]\arrow[dr,dash]\&{}\\
 {}\&R^*\cap V\arrow[dr,dash]\arrow[d,equal]\&{}\&\overline{(\varphi_0)}\arrow[dl,dash]\\
 {}\&\varphi_1\&\varnothing\&{}
\end{tikzcd}}
% L9–14, cancelled central diagram
\diagram{\begin{tikzcd}[ampersand replacement=\&,row sep=small,column sep=small,execute at end picture={\draw (\tikzcdmatrixname-1-2.north east)--(\tikzcdmatrixname-2-1.south west) (\tikzcdmatrixname-2-3.north east)--(\tikzcdmatrixname-3-2.south west) (\tikzcdmatrixname-2-3.north east)--(\tikzcdmatrixname-4-2.south west) (\tikzcdmatrixname-3-4.north east)--(\tikzcdmatrixname-4-3.south west);}]
 {}\&\deleted{R}\arrow[dl]\arrow[dr]\&{}\&{}\\
 \deleted{R^*}\arrow[dr]\&{}\&\deleted{\unsure{\Delta}=(\varphi_0,\varphi_1)}\arrow[dl]\arrow[dr]\&{}\\
 {}\&\deleted{R^*\cap\unsure{\Delta}=\varphi_1}\arrow[dr]\&{}\&\deleted{\varphi_0}\arrow[dl]\\
 {}\&{}\&\deleted{\varnothing}\&{}
\end{tikzcd}}
% L9–16, right diagram
\diagram{\begin{tikzcd}[ampersand replacement=\&,row sep=small,column sep=small]
 {}\&E\arrow[dl]\arrow[dr]\&{}\&{}\\
 \substack{\text{Objets}\\\text{couvrant }e_E}\arrow[dr]\&{}\&\substack{\text{sous-objets}\\\text{de }e_E}\arrow[dl]\arrow[dr]\&{}\\
 {}\&\substack{\text{objets}\\\text{finaux}}\arrow[dr]\&{}\&\text{objets vides}\arrow[dl,dash]\\
 {}\&{}\&\substack{\varnothing\ \text{sauf si }E=\mathrm{Top}(\varnothing)\\\text{où toutes les}\\\text{cat envisagées}\\\text{sont }=E=\\\text{cat ponctuelle}}\&{}
\end{tikzcd}}

% L13–17, left
\uline{NB} \deleted{\unsure{RR}} $V$ et $R^*\cap V$ n'ont pas de complémentaire mais ils ont les complémentaires faibles $\varnothing$ et $\varphi_0$

% L18–19, left
$V$\inserted{$\unsure{?}$} représente le \inserted{$?$} foncteur
\[
 E\longmapsto\unsure{\mathfrak P|\mathscr P|\mathcal P}(E)=\mathfrak P(e_E)
\]

% L16–21, right diagram
\diagram{\begin{tikzcd}[ampersand replacement=\&,row sep=small,column sep=small]
 {}\&\mathrm{Ens}\arrow[dl,dash]\arrow[dr,dash]\&{}\&{}\\
 \text{Ens. non vides}\arrow[dr,dash]\&{}\&\substack{\text{ens. de}\\\text{card }\leq 1}\arrow[dl,dash]\arrow[dr,dash]\&{}\\
 {}\&\substack{\text{ens.}\\\text{ponctuels}}\arrow[dr,dash]\&{}\&\{\varnothing\}\arrow[dl,dash]\\
 {}\&{}\&\varnothing\&{}
\end{tikzcd}}

% L21–23
\uline{NB}
\[
 \left\{\begin{array}{l}
 \text{Pas d'objets loc. constants, ou constants}\\
 \text{Pas d'objets \deleted{constants} qui \unsure{après un pr. sup|après une pr. sup} ont}\\
 \hspace{3em}\text{\deleted{\unsure{qui après}} (resp. id sur }e\text{)}\\
 \text{Pas d'objets constructibles}
 \end{array}\right.
\]

% L24–29, left (sideways)
% margin note written sideways
\margin{\unsure{C'est} donc des structures alg. définissables par $\underleftarrow{\lim}$ finies et même par produits finis, et sauf les deux derniers \unsure{$\Longleftarrow$} par \unsure{l'un des deux|l'un des trois} ---}

% L24–29
\uline{Question} Topos classifiant et sous-topos pour la structure
\[
 \begin{aligned}
 &\underbrace{X,Y}\quad\underbrace{f:X\to Y\ \text{sans plus}},\ \text{\unsure{ou|où}}\quad
 \underbrace{f:x-\unsure{x|X}\ \text{sans plus}}_{B_{\mathbb N^{\times}}},\quad \underbrace{f:e\to X}\\
 &\deleted{\unsure{\varphi}}\qquad X\times Y\to Z,\quad X\times Y\to X\\
 &\hspace{5em}X\times X\longrightarrow X\qquad !\\
 &\hspace{5em}X\times X\rightrightarrows x\qquad !!
 \end{aligned}
\]
% archive sheet number: 24
\scanpage{27}

% L1–2
\uline{Th.} Soit 1) $\sigma$ un ensemble de couples $(S,\lambda)$, $S$ catégorie, $\lambda=(\lambda',\lambda'')$ ens. de cônes (proj. et ind.) dans $S$

% L3–6
2) $\rho$ un ens de \deleted{diagrammes} \inserted{systèmes} $(\deleted{\unsure{R}},R_1,\delta_1,R_2,\delta_2,\varphi\deleted{\unsure{,R_3}})$ où $(R_i,\delta_i)$ \inserted{$(i=1,2)$} est une catégorie avec ens $\delta_i=(\delta'_i,\delta''_i)$ de cônes. $R$ une catégorie, \deleted{\unsure{$R_1,R_2$} des foncteurs \unsure{$\varphi_1$}} $\varphi$ un foncteur $R_2\mathrel{\overset{\varphi}{-}}R_1$ tel
% L5, right, cancelled triangle
\diagram{\begin{tikzcd}[ampersand replacement=\&,row sep=small,column sep=small,execute at end picture={\draw (\tikzcdmatrixname-2-1.south west)--(\tikzcdmatrixname-1-2.north east) (\tikzcdmatrixname-2-2.south west)--(\tikzcdmatrixname-1-3.north east);}]
{}\&\deleted{\unsure{R}}\arrow[dl,"\deleted{\varphi_1}"']\arrow[dr,"\deleted{\varphi_2}"]\&{}\\
\deleted{R_1}\&{}\&\deleted{R_2}
\end{tikzcd}}
que $\varphi(\delta_2)\subset\delta_1$ \deleted{$(R_1,\delta_2)$}

% L7
Soit $\Delta=\Delta_{\sigma,\rho}$ la \inserted{sous 2-}catégorie de Cat définie ainsi

% L8
a) Les objets sont les catégories $E$ ayant les propriétés suivantes
% L9–10
$a_1)$ $\forall(S,\lambda)\in\sigma$ et tout foncteur $S\xrightarrow{u}E$, les systèmes projectifs et inductifs correspondant aux cônes $\in u(\lambda)$ ont une limite
% L11–13
$a_2)$ $\forall(\deleted{\unsure{R}}R_1,\delta_1,\ldots)\in\rho$, \deleted{et tout foncteur $u_1:R_1-E$ qui est $\delta_1$-exact, il existe un foncteur $u_2:R_2-E$ qui est $\delta_2$-exact et tel que $u_1\varphi_1\simeq u_2\varphi_2$} \inserted{$u\mapsto u\circ\varphi$ induit une équivalence}
$\underline{\mathrm{Hom}}_{\delta_1}(R_1,E)\xrightarrow{\sim}\underline{\mathrm{Hom}}_{\delta_2}(R_2,E)$

% L14–19
b) Pour $E,E'\in\operatorname{Ob}\Delta$, $\underline{\mathrm{Hom}}_\Delta(E,E')$ est la sous-catégorie \inserted{(\unsure{str.|strte|stab.})} pleine de $\underline{\mathrm{Hom}}_{\mathrm{Cat}}(E,\unsure{E|E'|F})$ formée des foncteurs $F:E\to E'$ tels que l'on ait : $\forall(S,\lambda)\in\sigma$ et tout foncteur $u:S-E$ qui est $\lambda$-exact, $F\circ u$ est $\lambda$-exact. (Il est clair que cette condition est stable par isomorphisme et par composition, et vérifiée \deleted{par les} \inserted{si $F$ est une} équivalence de \unsure{catégorie|catégories}.)

% L20
On suppose satisfaites les conditions
% L21–25
A) $\sigma$ et $\rho$ sont petits \deleted{les} $\forall\,\unsure{(S,\lambda)|S,\lambda)}\in\sigma$ $S$ est petit et $\lambda$ est petit, et les \deleted{diagrammes} \inserted{catégories} d'indices des cônes $\in\lambda$ sont petites, de même pour tout $(\deleted{\unsure{R}},R_1,\delta_1,R_2,\delta_2,\varphi\deleted{\unsure{\varphi_2}})\in\rho$ les catégories \deleted{\unsure{$R_2$,}} $R_i$ $(i=1,2)$ sont petites, les $\delta_i$ $(i=1,2)$ sont \unsure{petit|petits} et les catégories d'indices des cônes $\in\delta_i$ sont petites.

% L26–27
B) $\forall(R_1,\ldots)\in\rho$, $(R_{\unsure{i|1}},\delta_{\unsure{i|1}})\in\sigma$ pour $i=1,2$. \inserted{(ou du moins : la condition d'exactitude des cat et \deleted{\illegible{1}} \unsure{des|de|le} foncteurs définie par $(R_i,\delta_i)$ est conséquence de celle définie par $\sigma$)} [On ne veut introduire des conditions d'exactitude supplémentaires b)
% archive sheet number: 26
\scanpage{28}

% L1–4
sur les catégories envisagées \inserted{$E$} pour $\Delta$ (en plus des conditions d'existence des lim dans a)) que si les \inserted{types de} lim que la formulation de ces conditions impliquent existent \uline{tjrs} dans les $E$ envisagées, et si les foncteurs admis entre ces $E$ y commutent ...]

% L5–7
Soit de plus $C$ une catégorie ess. petite, $d$ \inserted{$=(d',d'')$} un petit ensemble de \deleted{diag} \inserted{cônes} (proj. et ind.) \deleted{définis \unsure{p}} correspondant à des petites catégories d'indices.

% L8–11
Alors il existe un \uline{$C'\in\operatorname{Ob}\Delta$} et un foncteur $\boxed{i:C\unsure{-|\to}C'}$ qui est \uline{$d$-exact}, de telle façon que $(C',i)$ soit 2-\uline{univers}el, i.e. pour tout $E\in\operatorname{Ob}\Delta$, le foncteur $F\longmapsto F\circ i$ induit une \unsure{équiv|équiv.|équivalence}
\[
 \underline{\mathrm{Hom}}_{\Delta,i(d)}(C',E)\xrightarrow{\sim}\underline{\mathrm{Hom}}_{\mathrm{Cat},d}(C,E)
\]

% L12–13
De plus, cette catégorie $C'$ (définie à équivalence \deleted{uniq} près définie à isom. unique près) est ess. petite

% L14–15
\uline{Dém} On construit par réc. transfinie un \deleted{suite} \inserted{(ps) syst. inductif} transfinie $(C_i,d_i)$ de catégories sur $C$, munies d'un $d_i$ \inserted{$=(d'_i,d''_i)$} de cônes, en procédant ainsi
% L16–18
\[
 \begin{aligned}
 (C_0,d_0)&=(C,d)\\
 (C_{i+1},d_{i+1})&=K(C_i,d_i)\inserted{\quad\tau_i:C_i-C_{i+1}},\quad\text{où }K\text{ est une opération à expliciter plus bas}\\
 C_i&=\underset{j<i\,(\mathrm{Cat})}{\underrightarrow{\mathrm{Lim}}}\,C_j,\quad
 d_{\deleted{\unsure{j}}\unsure{i|j}}=\bigcup_{\underline j<i}(\underline{\mathrm I}\mathrm m\text{ de }d_j\text{ par }C_j-C_i)
 \quad\text{si }i\text{ ordinal limite}
 \end{aligned}
\]

% L19, interline
On note $\alpha_i:C\mathrel{\overset{\inserted{=C_0}}{-}}C_i$

% L20–22
La construction $K$ (indépendante de $i$) aura les propriétés suivantes, que pour la commodité des références nous écrivons avec les notations $(C_i,d_i),(C_{i+1},d_{i+1})$ :
% archive sheet number: 27
\scanpage{29}

% L1–2
a) $\forall E\in\operatorname{Ob}\Delta$, le foncteur $u\mapsto u\circ\tau_i$ induit une équivalence
\[
 \underline{\mathrm{Hom}}_{d_{i+1}}(C_{i+1},E)-\underline{\mathrm{Hom}}_{d_i}(C_i,E)
\]

% L3
b)
\[
 \tau_i(d_i)\subset d_{i+1}\subset\tau_i(d_i)\cup
 \bigcup_{(S,\lambda)\in\unsure{\rho|\sigma}}\ \bigcup_{u\in\underline{\mathrm{Hom}}_{\deleted{\unsure{\Delta}}}(S,C_{i+1})}u(\lambda)
\]

% L4
c) $\tau_i:C_{\unsure{i|i'}}-C_{i+1}$ est \uline{préexact} pour \deleted{\unsure{les}} $d_i$, i.e.

% L5–6, crossed block and arrow to e)
\deleted{\unsure{d)} $\forall(\unsure{R},R_1,\delta_1,R_2,\delta_2,\varphi\unsure{,\varphi_2})\in\rho$, et $u_1:R_1\to C_i$ qui est $\delta_1$-exact, \inserted{tel que $u_1(\delta_1)\subset d_i$,} il existe un foncteur $\delta_2$-exact $v_2:R_2\to C_{i+1}$ tel que $v_2\circ\varphi_2=(\tau_i u_1)\circ\varphi_1$ \inserted{(et $v_2(\delta_2)\subset d_{i+1}$)}}

% L7–12
d) $\forall(S,\lambda)\in\sigma$ \deleted{\unsure{et}} $u:S\longrightarrow C_i$, \deleted{et et une \unsure{diag|diag.} catégorie $I'$ $(I'')$ intervenant dans $\in\lambda'$ ou $\lambda''$} \inserted{$\unsure{\lambda}$ cône $\in\lambda'$ ou $\lambda''$} \deleted{(i.e. $\unsure{\lambda}=(\lambda',\lambda'')$), il existe un cône \inserted{$c'$} sur le syst. (proj ind) corresp. dans $\tau_i(c)$, tel que $c'\in d_{i+1}$} $\exists$ un foncteur $\bar v:\bar S\longrightarrow C_{i+1}$ (où $\beta:S-\bar S$ est obtenu en dédoublant les sommets des cônes critiques de $S$ $(**)$) tel que $\bar v\beta=\tau_i u$, et tel que $\bar v(\bar\lambda)\subset d_{i+1}$

% L10–12, slanted left
% slanted
\margin{cf définition sur p. $\rightarrow$ avant $(**)$}

% L13
\deleted{Prouvons que pour $\omega$ grand, $(C_\omega,d_\omega)$ est}

% L14–15
e) \quad \begin{math}\left\{\vcenter{\hbox{\begin{minipage}{.87\linewidth}
$e_0)$ $\forall(R_1,\delta_1,-)\in\rho$ et $u_1,v_1:R_1\rightrightarrows C_i$ \inserted{tels que $u_1(\delta_1)\subset d_i$, $v_1(\delta_1)\subset d_i$,} et $\alpha_1,\beta_1:u_1\rightrightarrows v_1$, la relation $\alpha*\varphi=\beta*\varphi$ implique $\tau_i*\alpha=\tau_i*\beta$
% L16–18
$e_1)$ $\forall(R_1,\delta_1,{-}{-})\in\rho$ et $u_1,v_1:R_1\rightrightarrows C_i$ tels que $u_1(\delta_1),v_1(\delta_1)\subset d_i$, et tout hom $u_1\circ\varphi\xrightarrow{\alpha_2}v_1\circ\varphi$, $\exists$ un hom $\tau_i u_1\xrightarrow{\bar\alpha_1}\tau_i v_1$ tel que \deleted{\unsure{e)}} $\bar\alpha_1*\varphi=\tau_i*\alpha_2$
% L19–20
$e_2)$ $\forall(R_1,\delta_1,{-}{-})\in\rho$ et $u_2:R_2-C_i$ tel que $u_{\unsure{2|1}}(\delta_2)\subset d_i$, $\exists\,\bar u_1:R_1\longrightarrow C_{i+1}$ tels que $\bar u_1\circ\varphi=\tau_i\circ u_2$ et $\bar u_1(\delta_1)\subset d_{i+1}$
\end{minipage}}}\right.\end{math}
% archive sheet number: 28
\scanpage{30}

% L1–2
De a) on déduit, par récurrence transfinie, pour tt $E\in\operatorname{Ob}\Delta$
\[
 \underline{\mathrm{Hom}}_{\deleted{\unsure{\Delta}}d_i\deleted{\unsure{\rho}}}(C_i,E)\xrightarrow{\sim}
 \underline{\mathrm{Hom}}_d(C,E)\qquad(F\mapsto F\circ\alpha_i)\tag*{(1)}
\]

% L3–4
De b) on déduit de $\overline{\mathrm m}$
\[
 \alpha_i(d)\subset d_i\unsure{\subset|=}\alpha_i(d)\cup
 \bigcup_{\substack{\underline{j}<i\\(S,\lambda)\in\unsure{\rho|\sigma}}}
 \bigcup_{u\in\underline{\mathrm{Hom}}(S,C_j)}
 \tau_{ij}u(\lambda)\qquad\substack{\text{si }i\text{ ordinal}\\\text{limite}}\tag*{(2)}
\]

% L5–6
Prouvons que pour $\omega$ \inserted{ordinal limite} assez grand ($\omega$ un ``petit'' ordinal), $(C_\omega,d_\omega)$ est une solution du Pb. \deleted{Cela signifie} \inserted{Il suffit} qu'on ait
% L7–9
$1^\circ)$ $C_\omega\in\operatorname{Ob}\Delta$ i.e. \quad $a_1)$ $\forall(S,\lambda)\in\sigma$ et $u:S-C_\omega$, existence des lim correspondant aux cônes $\in u(\lambda)$
\par $a_2)$ $\forall(\unsure{R_0},R_1,\ldots\varphi)\in\rho$ \deleted{et $u_1:R_1\to C_\omega$ \inserted{$\delta_1$-exact}, $\exists\,u_2:R_2\to C_\omega$ \inserted{$\delta_2$-exact} et $u_1\varphi_1\simeq u_2\varphi_2$}
\[
 \begin{aligned}
 &\underline{\mathrm{Hom}}_{\delta_1}(R_1,C_\omega)\\[-2pt]
 &\hspace{2em}\unsure{\wr|\simeq}\\[-2pt]
 &\underline{\mathrm{Hom}}_{\delta_2}(R_2,C_\omega)
 \end{aligned}
\]

% L10
$2^\circ)$ $\alpha_{\unsure{i|\omega}}:C\longrightarrow C_\omega$ est $d$ exact

% L11
$3^\circ)$ \deleted{\unsure{$E$} \illegible{1}}
\[
 d_\omega=\alpha_\omega(d)\cup\bigcup_{(S,\lambda)\in\sigma}
 \bigcup_{u\in\underline{\mathrm{Hom}}_\lambda(S,C_\omega)}u(\lambda)
\]

% L12–13
[Car $3^\circ$ implique que pour tte $E\in\operatorname{Ob}\Delta$ on a
\[
 \underline{\mathrm{Hom}}_{\Delta,\alpha_\omega(d)}(C_\omega,E)\deleted{\unsure{\simeq}}=
 \underline{\mathrm{Hom}}_{\deleted{\unsure{\Delta,\alpha_\omega(d)}}d_\omega}(C_\omega,E)
 \xrightarrow[(1)]{\sim}\underline{\mathrm{Hom}}_d(C,E)
\]

% L14, interline
\deleted{\unsure{donc|Alors}} Soit $T_\omega$ l'ensemble des ordinaux $j<\omega$

% L15–16
Pour vérifier $2^\circ)$ il suffit de vérifier l'\unsure{assert.|assertion} plus forte
\[
 \forall i<\omega,\quad\tau_{\omega,i}:C_i-C_\omega\text{ est }d_i\text{-exact}\tag*{(3)}
\]

% L17–23
Ceci résulte à son tour de c) Si $T_\omega$ est grand devant les cardinaux : $\operatorname{card}\operatorname{Fl} I$, où $I$ est une cat. d'indices intervenant dans un des cônes $\in d_i$ ; Or en vertu de (2), ces $I$ \deleted{sont \unsure{toutes}} sont \unsure{ou bien|au plus|de plus} ceux qui interviennent dans $d$, soit dans un $\lambda$ associé à un $(S,\lambda)\in\sigma$ -- c'est un petit ens de petites catégories, donc les $\operatorname{card}\operatorname{Fl}I$ sont majorés par un cardinal $\pi$ -- et il suffit $T_\omega$ grand devant $\pi$.

% L18–19, left (slanted)
% slanted
\margin{\unsure{dégager|désigner} lemme}
% archive sheet number: 29
\scanpage{31}

% L1–2
Il résulte de (3) et (2) appliqué à $i=\omega$ que l'on a
\[
 \alpha_\omega(d)\subset d_\omega\subset\alpha_\omega(d)\cup
 \bigcup_{\substack{\deleted{\unsure{i<\omega}}\\(S,\lambda)\in\sigma}}
 \ \bigcup_{u\in\underline{\mathrm{Hom}}_\lambda(S,C_\omega)}
 \qquad\substack{\text{mod isom de}\\\text{cônes}}
\]
% L3–5
et pour prouver $3^\circ)$ il faut donc prouver
\[
 \begin{aligned}
 &\forall(S,\lambda)\in\sigma\text{ et }u\in\underline{\mathrm{Hom}}_\lambda(S,C_\omega),
 \ \deleted{\text{exi}}\exists\,i<\omega\text{ et un}\\
 &\text{hom }u_i:S\to C_i\text{ tel que }u_i(\lambda)\subset d_i\ \inserted{\text{(mod iso de cônes)}}
 \text{ et }u\simeq\tau_{\omega,i}\circ u_i
 \end{aligned}\tag*{(\deleted{\unsure{3 bis}}\inserted{4})}
\]

% L6–7
\deleted{Cela résulte facilement du fait que tt se ramène à un $C_i$, de d) et du fait que}
Enfin, \deleted{\unsure{pas}} examinons $1^\circ)$. \deleted{A tout \unsure{prendre}}

% L8–10
$a_1)$ Sera conséquence de d) si nous notons que tout foncteur $u:S-C_i$ se factorise par un $C_i$ -- pourvu qu'on prenne $\pi\geq\operatorname{Card}\operatorname{Fl}S$ pour tt $(S,\lambda)\in\sigma$.

% L11–12
$a_2)$ Sera conséquence de e), utilisant B (les $(R_i,\delta_i)\in\sigma$)) et \deleted{\unsure{(3 4 i)|(4)}} et (3) (pour $i+1$ au lieu de $i$).

% L13–14, crossed block
\deleted{Donc il reste à \unsure{donc|donner} dégager des conditions sur la construction \unsure{$\alpha|K$} qui assurent $(\unsure{4|3})$, et d' \inserted{enfin à} expliciter une construction $K$.}

% L15, left
\margin{peut remonter !}

% L15–18
Prouvons 4). On sait que $u:S-C_\omega$ se factorise par un \inserted{$u_i:S\to$} $C_i$ (pourvu qu'on prenne $\pi\geq\operatorname{card}\operatorname{Fl}S$ $\forall(S,\lambda)\in\sigma$), donc quitte à composer avec $\tau_i$, on peut supposer que \unsure{cel|cela|ceci} se prolonge en $\bar u_i:\bar S-C_i$ tel que $\bar u(\bar\lambda)\in d_i$, donc $\tau_{\omega,i}\circ\bar u_i$ est un hom $\bar S\to C$ qui est $\bar\lambda$ exact. \inserted{grâce à (3).} \deleted{Il est donc} Comme
% L19–24
$u$ est $\lambda$-exact, $\bar u(i_c)$ est un isom pour tt $c\in\lambda$, donc $\bar u_i(i_c)$ devient un isom pour tt $c\in\lambda$, donc les $\tau_{ji}\bar u_i(i_c)=\bar u_j(i_c)$ ont la $=$ propriété pour $j\in T_\omega$ assez grand (pourvu que $\pi\geq\operatorname{card}\lambda=\operatorname{card}\{c\}$) donc \deleted{\illegible{1}} les \deleted{éléments} cônes de $u_j(\lambda)$ sont \uline{isom} à ceux de $\bar u_j(\bar\lambda)$, comme ceux-ci sont dans $d_i$, il en est de $=$ de ceux là à isom près, cqfd.
% archive sheet number: 30
\scanpage{32}

% L0, upper left (slanted)
% slanted
\margin{$*$ \uline{page \unsure{annexe|arrière}}}

% L0, upper right, detached marks
\margin{$\cdot\!\cdot\quad\cdot\ \cdot$}

% L1–5
$\bigcirc\!\!\!{\scriptstyle\genfrac{}{}{0pt}{}{*}{*}}$ Soit $S$ une catégorie avec $\lambda=(\lambda',\lambda'')$ ens de cônes. On lui associe une catégorie $\bar S$ et un foncteur $S\xrightarrow{\beta}\bar S$ \deleted{tel que} \deleted{\unsure{\(\beta\)|p}} en ajoutant, pour tt \deleted{sommet de} cône \inserted{$c':(x\xrightarrow{p_\alpha}(x_\alpha))$} $\in\lambda'$ \deleted{\unsure{$x\to$}}, un objet $\bar x_c$ à $S$, avec \deleted{une de} \inserted{une flèche $x\to\bar x_c$ et} une famille de flèches $\bar x_c\xrightarrow{q_\alpha}x_\alpha$ factorisant les $p_\alpha$ -- et de même dualmt pour les cônes $c''\in\lambda''$.
% L4–7, left
\margin{\diagram{\begin{tikzcd}[ampersand replacement=\&,row sep=small,column sep=small]
 {}\&{}\&x_1\\
 x\arrow[urr]\arrow[r,"i_c"']\arrow[rr]\arrow[drr]\&\bar x\arrow[ur]\arrow[r]\arrow[dr]\&x_2\\
 {}\&{}\&x_3
\end{tikzcd}}}

% L6–9
On dit que $(\bar S,\bar\lambda)$ est déduit de $(S,\lambda)$ en \uline{dédoublant} les sommets des cônes critiques de $S$ ; $\bar\lambda$ est formé des cônes déduits des cônes $c\in\lambda$ en remplaçant le sommet $x_c$ par $\bar x_c$.

% L9–12
Donc pour tt foncteur $u:S-\unsure{\bar E|E}$, avec $E\in\operatorname{Ob}\Delta$, $\exists$ une \uline{unique} extension en $\bar u:\bar S-E$ qui soit $\bar\lambda$-exact, de plus, $u$ est $\lambda$-exact ssi \deleted{\unsure{tt les|ses}} pour tt $c\in\lambda$, $\bar u(i_c)$ est un isomorphisme.

% L13–14
$(*)$ Définition de foncteurs \uline{$\delta$-préexacts} $D.\xrightarrow{u}F$, $D$ muni d'un ens $\delta=(\delta',\delta'')$ de cônes....
% archive sheet number: 31
\scanpage{34}

% L1–2, across top
\[
 H^n(K(G,1),\unsure{\mathbb Z|\Pi})\qquad
 \overset{\unsure{3}}{\boxed{\underline{H^n(\underline{K(G,1)}\unsure{,|;}\Pi)}}}
 \underset{Z^{\unsure{n|i}}(\Pi)}{\overset{\unsure{\sim|n}}{\longrightarrow}}\mathrm{Simpl}(\mathrm{ENS})
\]
\deleted{\unsure{$H$}}

% L3–4
Lawvere theory \qquad $\mathcal L$ a small cat. \qquad $1\in\mathcal L$ \qquad $n=\amalg\,1$
\[
 (\mathcal L^o\to\mathrm{ENS})\qquad \mathrm{Mod}(\mathcal L^o)
\]

% L4–7, upper left
\diagram{\begin{tikzcd}[ampersand replacement=\&,column sep=large,row sep=large]
 S=\underline{\mathrm{Hom}}^{\prime}(C^o,\mathrm{ENS})\arrow[r,leftrightarrow,"\sim"]\arrow[d,"U"]\&
 \boxed{\underline{\mathrm{Alg}.(\unsure{\mathbb T|T|\mathfrak T})}}\arrow[l,bend right=22,dashed]\arrow[dl,"U"]\arrow[d]\\
 (\mathrm{ENS})\arrow[u,bend left,"F"]\arrow[ur,bend left=15]\arrow[r,"\sim"']\&{}
\end{tikzcd}}

% L5–10, upper right
$P\in\operatorname{Ob}(S)$

% L6, upper right, first condition
$1^\circ$\quad $P_{\underline I}=\amalg P\qquad P\times\unsure{\mathbb I|I}.$

% L7, upper right, second condition
$2^\circ$\quad Every eq. rel is effective \qquad $R\rightrightarrows S$

% L8, upper right, third condition
$3^\circ$\quad \uline{$f:X\to Y$ effective $\mathrel{\overset{\text{\unsure{if 2)|epi}}}{\Longleftrightarrow}}\mathrm{Hom}(P,X)\xrightarrow{\text{surjective}}\mathrm{Hom}(P,Y)$.}

% L9, upper right, fourth condition
$4.$\quad \uline{fiber prod} \unsure{?}

% L10, upper right, fifth condition
$5^\circ$
\diagram{\begin{tikzcd}[ampersand replacement=\&,row sep=small]
 P\arrow[r]\arrow[dr,"\substack{n\\\mathrm{fin}}"']\&P\times K\\
 {}\&P\times n\arrow[u,hook]
\end{tikzcd}}

% L9–10, beside list
% slanted
\margin{$\searrow$ \unsure{Haks} $\underline I$\quad $+$ \quad $\times$}

% L9, left
\diagram{\begin{tikzcd}[ampersand replacement=\&,column sep=small,row sep=small]
 G_1\arrow[rr]\arrow[dr,"\unsure{2}"']\&{}\&G_2\arrow[r,shift left]\arrow[r,shift right]\&
 |[draw,rounded corners]|{\unsure{G_{\dots}*G_{\dots}|S*S_1}}\\
 {}\&\unsure{C}\arrow[ur]\&{}\&{}
\end{tikzcd}}

% L11–14, left
\diagram{\begin{tikzcd}[ampersand replacement=\&,row sep=small]
 {}\arrow[r]\&(\unsure{/})\arrow[r]\arrow[d,shift left]\arrow[d,shift right]\&R\arrow[d,shift left]\arrow[d,shift right]\\
 {}\&\unsure{T'|T}\arrow[r,"\unsure{h}"]\arrow[d,dash]\&X\arrow[d,"\unsure{=}"']\\
 {}\&T\arrow[r]\&X
\end{tikzcd}}

% L11–14, centre
$P$ \quad morphism of \unsure{de|\deleted{the}} \uline{effective} \uline{descent}
\diagram{\begin{tikzcd}[ampersand replacement=\&,column sep=small,row sep=small]
 {}\&\mathrm{Hom}(P,X)\arrow[dl]\\ +\&X\arrow[r,"f"']\&Y
\end{tikzcd}}
\unsure{\uline{Vergl} pre \unsure{pp.}}

% L13–18, centre right
\diagram{\begin{tikzcd}[ampersand replacement=\&,row sep=large,column sep=large]
 \mathbb A\arrow[r,"\Phi","\sim"']\arrow[d,"U"]\&
 |[draw,rounded corners,double]|{\mathrm{Alg}(\mathbb T)}\arrow[l,bend right=48]\arrow[dl]\\
 \underline{\mathbb B}\arrow[u,bend left,"F"]\arrow[ur,bend right]\arrow[r,"T","\unsure{\mathbb T}"']\&T\mathbb B
\end{tikzcd}}

% L13–15, right
\margin{\unsure{Ben|Beck} \unsure{'Barr|'Baur|Barr} \uline{Linton} \qquad $U$ tripleable}

% L15–17, left and centre
$\mathrm{ENS}\longrightarrow\mathcal L$
\diagram{\begin{tikzcd}[ampersand replacement=\&,row sep=small]
 C/Y\arrow[d,"f^*"]\\ C/X\arrow[u,bend left,"f_!"]\quad\unsure{(}
\end{tikzcd}}

% L16–19, right
\diagram{\begin{tikzcd}[ampersand replacement=\&,row sep=small,column sep=small]
 T\arrow[r]\arrow[d,shift left]\arrow[d,shift right]\&R\arrow[d,shift left]\arrow[d,shift right]\\
 T_1\arrow[r]\arrow[d,"\unsure{+}"']\&X\arrow[d]\\ Q\arrow[r]\&Y
\end{tikzcd}}

% L16–19, left and middle
\diagram{\begin{tikzcd}[ampersand replacement=\&,row sep=small,column sep=small]
 {}\&T(\&{}-U(G_2)\\
 {}\arrow[r,shift left]\arrow[r,shift right]\&\underline{U(FUG)}\arrow[r,"\varepsilon"]\&U(G)\arrow[u,"U(f)"]\\
 {}\&\mathbb T(UG)\arrow[r]\&(\unsure{U(-)|UG})
\end{tikzcd}}
$G\qquad\underline G\qquad\mathbb T=(\underset{\shortparallel}{\overset{\inserted{UF}}{\deleted{\unsure{VF}}}},-,-)$ \qquad $T$

% L20–21, left
\diagram{\begin{tikzcd}[ampersand replacement=\&,column sep=small]
 \unsure{T}\arrow[r,shift left]\arrow[r,shift right]\&\mathbb T(X)\arrow[r,"f"]\&X
\end{tikzcd}}
$\unsure{\mathrm{As}|A}\rightleftarrows\mathrm{ENS}$ \qquad $\unsure{\mathfrak d}\ u$ \qquad $UF(G)\longrightarrow UG$

% L20–23, right
$U$ \uline{\unsure{trip.}}, \qquad \uline{\unsure{fin}} \quad $|$

% L22, right, numbered conditions
$\ldots$ \uline{iff} ; \qquad $1^\circ$ $U$ is \deleted{conservative} \inserted{\uline{\unsure{conserv|conserve}}} \qquad $2^\circ$

% L23–26, left
\diagram{\begin{tikzcd}[ampersand replacement=\&,column sep=small,row sep=small]
 {}\&{}\&T\&{}\\
 UF(\quad)\arrow[rr]\&{}\&(\ )\&{}\\
 \unsure{F^2(S)}\arrow[r,shift left,"F(f)"]\arrow[r,shift right,"\mu"']\&F(\overset{\unsure{Xid}}S)\arrow[r,"f"]\&S\arrow[l,bend left,"\eta"]\arrow[uu,"x"]\arrow[r,mapsto]\&|[draw,circle]|S\\
 (UF)(UF)(X)\arrow[rr]\&{}\&X\&X\in\operatorname{Ob}(\mathbb B)
\end{tikzcd}}

% L26–27, left
\diagram{\begin{tikzcd}[ampersand replacement=\&,column sep=small]
 {}\arrow[r,shift left]\arrow[r,shift right]\&T^2(\quad)\arrow[r,"\unsure{F^2}"]\&{}\arrow[r,bend right]\&F(G)\arrow[r]\&G
\end{tikzcd}}

% L25–30, right
\diagram{\begin{tikzcd}[ampersand replacement=\&,row sep=small,column sep=small]
 X_1\arrow[r,shift left]\arrow[r,shift right]\&X_{\unsure{0|i|2}}\arrow[r]\&Q'\ \text{in }\mathbb A\\
 {}\&{}\&U\\
 {}\&{}\&U(Q')\\
 U(X_1)\arrow[r,shift left]\arrow[r,shift right]\&U(X_{0})\arrow[l,bend left]\arrow[r]\&Q\arrow[l,bend left]\arrow[u,"\unsure{\iota}"]
\end{tikzcd}}

% L28–30, left
\diagram{\begin{tikzcd}[ampersand replacement=\&,column sep=small]
 T^2\arrow[r,shift left,"T(f)"]\arrow[r,shift right,"\mu"']\&TX\arrow[r,"f"]\&X\arrow[l,bend left,"\eta"']\ \mathrm{id}
\end{tikzcd}}
\[
 T(X)\mathrel{\overset{\xi}{-}}/X
\]

% L31–33, foot
\[
 \begin{aligned}
 1^\circ\quad p_0s_1&=s_0p_{\unsure{0|1}}\\
 2^\circ\quad p_1s_1&=\mathrm{id}\\
 p_0s_0&=\mathrm{id}
 \end{aligned}
\]
\diagram{\begin{tikzcd}[ampersand replacement=\&,row sep=small,column sep=small]
 X_1\arrow[r,shift left,"p_0"]\arrow[r,shift right,"p_1"']\& X_0\arrow[l,bend right=40,"s_1"']\arrow[r,"p_0"']\&X_{-1}\arrow[l,bend right=40,"s_0"']\\
 {}\&+_{-1}\arrow[u,shift left,"{\unsure{s_1p_0}}"']\arrow[u,shift right,"\mathrm{id}"]\&{}
\end{tikzcd}}
$\displaystyle s_1=(\unsure{s_1p_0},\mathrm{id}_{X_0})\qquad\Longrightarrow\quad X'\rightrightarrows X\rightrightarrows Y$
% archive sheet number: 33
\scanpage{35}

% L1 top centre, cancelled
\deleted{\unsure{Exem|théo}}

% L2 top left, pencil
\diagram{\begin{tikzcd}[ampersand replacement=\&,execute at end picture={\draw[thin] ([xshift=-.12cm]\tikzcdmatrixname-1-1.north west) to[bend left=42] (\tikzcdmatrixname-2-1.south east) to[bend left=48] cycle;}]
X_{\unsure{i|1}}' \arrow[d] \& U X_{\unsure{0|i}}\arrow[d]\\
X \& U X\arrow[u,bend right=40]
\end{tikzcd}}

% L3 upper left and right
$C/\unsure{S|\zeta}$
\diagram{\begin{tikzcd}[ampersand replacement=\&]
X_1\arrow[r,shift left=.5ex,"x_1"]\arrow[r,shift right=.5ex,"x_0"'] \& X_0\arrow[r,leftrightarrow] \& u
\end{tikzcd}}
\diagram{\begin{tikzcd}[ampersand replacement=\&]
u\quad\deleted{\unsure{u|w}}\quad\deleted{\unsure{a}} \& {}\arrow[d] \& x\arrow[d,mapsto]\\
\& {} \& X'
\end{tikzcd}}
$\unsure{G|C}/{\unsure{s|S}}'$

% L4 upper right, bracket and split diagram
\diagram{\begin{tikzcd}[ampersand replacement=\&,execute at end picture={\draw (\tikzcdmatrixname-1-1.north west) -- (\tikzcdmatrixname-1-2.north east) to[bend right] ([yshift=-2pt]\tikzcdmatrixname-1-2.north east);}]
\underline{X_1}\arrow[r,shift left=.7ex,"d_0"]\arrow[r,shift right=.65ex,"\deleted{d_1}"']\arrow[d] \&
\underline{X_{1}}\arrow[l,bend left=35,"s_1"']\arrow[r,"p"]\arrow[d] \&
\underline Q\arrow[l,bend left=30,"s_0"']\\
Q\arrow[u,bend left=35]\arrow[r,shift left=.5ex]\arrow[r,shift right=.5ex] \& Q\arrow[u,bend left=30] \&
\end{tikzcd}}
\[
\left\{\begin{aligned}
d_1s_1&=\mathrm{id}_{X_1}\\
d_0s_1&=s_0p
\end{aligned}\right.
\]

% L5 left, slanted heading and triangles
% slanted
\margin{$I$ espace discret}
\diagram{\begin{tikzcd}[ampersand replacement=\&]
\& (\mathrm{Ens})^I\arrow[d,"\wr" description,draw=none]\\
B_\pi\arrow[dr] \& B=\mathrm{Top}(\unsure{E|\Delta|I})\arrow[l]\arrow[d,"f"]\\
\& B'
\end{tikzcd}\qquad
\begin{tikzcd}[ampersand replacement=\&]
 B\arrow[d,shift right=.65ex,"f_*"']\&Y\\
 B'\arrow[u,shift right=.65ex,"f^*"',very thick]\&X
\end{tikzcd}}

% L6 centre, endofunctor and two maps
\[
\pi=f^*f_*\colon B-B\quad \boxed{\underline{\text{exact à g.}}}\quad\uline{\text{accessible}}
\qquad
\left\{\begin{aligned}\pi&-\mathrm{id}\\ \pi&-\pi^2\end{aligned}\right.
\]
\diagram{\begin{tikzcd}[ampersand replacement=\&]
\deleted{\unsure{\text{Ens}|B}}\arrow[r] \& \underline{\mathrm{coAlg}(B,\pi)}\arrow[r,very thick]\& B\\
\&\underline{\text{topos}}=B_\pi\&
\end{tikzcd}\qquad
\begin{tikzcd}[ampersand replacement=\&]B\arrow[d,"\deleted{\text{fidèle}}\ \inserted{\text{conservatif}}"]\\ B_\pi\end{tikzcd}}

% L7 middle, criteria
\[
\underset{\text{tous deux}}{[\text{conserv.tif}]}\ 
\overset{\mathrm{déf}}{\Longleftrightarrow} f^*\ \text{fidèle}\ (= \underline{\text{conserv.tif}})\
\overset{\mathrm{th}}{\Longleftrightarrow}\underline{B_\pi\xrightarrow{\sim}B}
\]
\[
\left\{\begin{aligned}
& f^*\ \text{conservatif}\\
& f^{\unsure{\times|*}}\ \text{commute aux noyaux de couples}
\end{aligned}\right\}
\]

% L8 lower middle, adjunction and overwritten term
\[
\begin{aligned}
\mathrm{Hom}_{\mathrm{Ens}}(f^*_{\unsure{G}}X,Y)&\simeq\mathrm{Hom}_G(X,f_*(Y))\qquad
f_*(Y)=\underline{\mathrm{Hom}(G,Y)}\\
&\hspace{4em}\wr\\[-.7ex]
&\underline{\mathrm{Hom}}_{\unsure{G|\mathrm{Ens}}}\,(X,\deleted{f_*(Y)}\mathrm{Hom}_{\mathrm{Ens}}(G,Y))
\qquad Y-\mathrm{Hom}(G,Y)
\end{aligned}
\]

% L9 lower left, existence question and paired diagrams
Soit $B$ un topos

$\exists\ B'$ topos avec assez de pts
\diagram{\begin{tikzcd}[ampersand replacement=\&]
\widetilde C\arrow[r,"\simeq" description,draw=none]\arrow[d,"\unsure{\rho|\hookrightarrow}" description,draw=none]\& B\arrow[d,"f"] \&
B\arrow[d,shift right=.6ex,"f_*"']\\
\widehat C\& B'\&B'\arrow[u,shift right=.6ex,"f^*"']
\end{tikzcd}}

% L10 lower right, group calculations
\diagram{\begin{tikzcd}[ampersand replacement=\&]
\pi(Y)\arrow[d]\arrow[r,"=" description,draw=none] \& \mathrm{Hom}(G,Y)\arrow[d] \& e\arrow[d]\\
Y\&Y\&G\\
\&\&G^{\unsure{2|1}}-G
\end{tikzcd}}
\[
\begin{aligned}
\pi&-\pi^2\\
\mathrm{Hom}(G,Y)&-\mathrm{Hom}(G,\mathrm{Hom}(G,Y))\\
&\hspace{7em}\wr\\[-.7ex]
&\hspace{5em}\mathrm{Hom}(G^{2},Y)\\
&G\leftarrow G^{2}
\end{aligned}
\]

% L11 bottom, enclosing rule
\figure{\begin{tikzpicture}[x=1cm,y=1cm]
\draw (0,.75) .. controls (.75,.62) and (3.7,.99) .. (5.8,1.00) .. controls (6.0,1.00) and (5.82,.10) .. (6.22,.09) .. controls (6.75,.04) and (7.82,.06) .. (8.60,.09);
\draw (8.90,.10)--(9.15,.10);
\end{tikzpicture}}
% archive sheet number: 34
